\documentclass[10pt,twocolumn]{article}
\usepackage[utf8]{inputenc}
\usepackage[english,russian]{babel}
\usepackage{amsmath,amssymb,mathtools,bm,geometry,amsthm,booktabs,tikz-cd}
\geometry{a4paper,margin=1.5cm}

\DeclareMathOperator{\Rank}{Rank}
\DeclareMathOperator{\Coker}{Coker}
\DeclareMathOperator{\Tr}{Tr}
\DeclareMathOperator{\Aut}{Aut}
\DeclareMathOperator{\Hom}{Hom}
\DeclareMathOperator{\Spec}{Spec}
\DeclareMathOperator{\Gal}{Gal}

\newtheorem{law}{Закон}
\newtheorem{criticism}{Замечание}

\begin{document}

\title{\textbf{Арифметико-геометрические основания мультимасштабных NoC-структур: Проконечные пополнения, мотивы и локально-глобальный принцип следа}}
\author{\textbf{Унифицированная комбинаторно-алгебраическая модель}}
\date{Сентябрь 2026}
\maketitle

\section*{Введение и категорный базис}
Динамическое распределение трафика пакетов в коммутационных каскадах сетей на кристалле (NoC), формализованное через соотношение симметрической группы $S_n$ и пространства беспорядков $D_n$, допускает интерпретацию в терминах оснований математики. Мы рассматриваем асимптотическую эволюцию топологии NoC как действие на проконечном пополнении кольца целых чисел:
\begin{equation}
\hat{\mathbb{Z}} = \varprojlim \mathbb{Z}/n\mathbb{Z} \cong \prod_{p} \mathbb{Z}_p
\end{equation}
Переход от локальных дефектов к глобальному объему задается семейством функторов из категории конечных групп $\mathbf{Grp}$ в категорию арифметических схем $\mathbf{Sch}_{/\mathbb{Q}}$.

\subsection*{1. Башня мотивных расширений и мультимасштабный синтез}
Пусть $X$ — арифметическая схема, ассоциированная с регулярным направленным графом $\mathcal{G}$, индуцированным циклическим разностным множеством (ЦРМ) $D$ с параметрами $(v=157, k=13, \lambda=1)$ над конечным полем $\mathbb{F}_v$. Базовые этапы деградации топологии фиксируются через иерархический базис порядков проективных плоскостей $N(d) = d^2 + d + 1$:
\begin{equation}
\{N(3), N(6), N(10), N(12)\} = \{13, 43, 111, 157\}
\end{equation}
Полная комбинаторная емкость конфигурации изоморфна порядку знакопеременной группы $\vert A_6 \vert = 360$ и разлагается в прямую сумму мотивов Гротендика, управляемых орбитами симметрии подгруппы $S_4$:
\begin{equation}
\vert A_6 \vert = N(12) + N(10) + N(6) + N(3) + \frac{3}{2}\vert S_4 \vert = 360
\end{equation}

Редукция критических сингулярностей $\Phi_6 = 43$ и $\Phi_{10} = 111$ по модулю локальной связности $k=13$ индуцирует коядро, кодирующее размерность пространства гомотопических препятствий (аналог кручения в группах Селмера):
\begin{equation}
\Coker\left([\Phi_{10}] \ominus [\Phi_{6}]\right) \equiv (111 - 43) \pmod{13} \equiv 3
\end{equation}

В мультирезонансных составных базах $B \in \{6, 12, 60\}$ диофантова разность проецируется на жесткую сетку делителей кольца аделей $\mathbb{A}_{\mathbb{Q}}$, образуя зазор:
\begin{equation}
V_{60} \ominus (\Phi_{10} \oplus \Phi_6)_{60} = (2.37)_{60} \ominus (2.34)_{60} = (0.03)_{60} \equiv 3_{10}
\end{equation}

Мультимасштабная интеграция канонических проекций отображается диаграммой спуска:
\[
\begin{tikzcd}
\mathbb{A}_{\mathbb{Q}} \arrow[r, "\text{Res}_{B=6}"] \arrow[d, "\text{Res}_{B=60}"'] & (111)_6 \arrow[d, "\Coker"] \\
(0.03)_{60} \arrow[r, "\cong"] & 3_{10} \pmod{13}
\end{tikzcd}
\]

\subsection*{2. Фундаментальные законы арифметики NoC}

\begin{law}[О разрядно-инверсной дуальности в $\hat{\mathbb{Z}}$]
В факториальной системе счисления (микст-базисе) $X = \sum_{i=1}^{n-1} c_i \cdot i!$, коэффициенты субфакториала $!n$ и пространства фиксированных точек $\Delta_n = n! - !n$ при $n=7$ удовлетворяют условию инверсного дополнения с точностью до циклического переноса в $\mathbb{Z}$-модуле:
\begin{equation}
c_i(!7) + c_i(\Delta_7) = i, \quad \forall i \in \{3, 4, 5, 6\}
\end{equation}
\end{law}

\begin{proof}
Вычисление факториальных кодов для $\!7 = 1854$ и $\Delta_7 = 3186$ остатками на $i!$ дает точные векторы: $\!7 = (232100)_{!}$ и $\Delta_7 = (422300)_{!}$. Поразрядная сумма весов $c_i(\!7) + c_i(\Delta_7)$ дает $6, 5, 4$ для старших разрядов, фиксируя инвариантную меру Якоби на проконечных многообразиях.
\end{proof}

\begin{law}[О трансцендентном регуляторе знакопеременных мотивов]
Полуразность $\Xi_n = |A_n| - !n = \Delta_n - |A_n|$ выступает тождественным арифметическим регулятором. При $n \to \infty$ нормированный сдвиг $\Xi_{\text{rel}}$ сходится к трансцендентной константе:
\begin{equation}
\Xi_{\text{rel}} = \lim_{n \to \infty} \frac{\Xi_n}{n!} = \frac{1}{2} - e^{-1} \approx 0.132120558828\dots
\end{equation}
при этом бесконечные мантиссы дробных частей $\{ \Xi_{\text{rel}} \}$ и $\{ P_{\text{ord}} \} = \{ 1 - e^{-1} \}$ поразрядно изоморфны с трансляционным сдвигом на рациональную константу $0.5$.
\end{law}

\begin{law}[О регулярном blow-up сингулярных дефектов]
Локализация критического дефекта $\Phi_6 = 43$ в кольце вычетов по составному модулю $B=6$ эквивалентна операции раздутия (blow-up) особенности арифметической схемы в гладкий дивизор, представленный автоморфным монолитом:
\begin{equation}
43_{10} = 1 \cdot 6^2 + 1 \cdot 6^1 + 1 \cdot 6^0 = (111)_6
\end{equation}
\end{law}

\begin{law}[О локально-глобальном принципе следа]
Динамический шаг комбинаторного взрыва $\Delta_{13} = 13! - !13$ удовлетворяет жестким ограничениям жесткости схемы над ортогональными модулями связности ($k=13$) и глобального объема ($V=157$):
\begin{equation}
\Delta_{13} \equiv 1 \pmod{13} \quad \text{и} \quad \Delta_{13} \equiv 13 \pmod{157}
\end{equation}
\end{law}

\begin{proof}
По теореме Вильсона $13! \equiv 0 \pmod{13}$. Из свойств субфакториала следует $\!13 \equiv -1 \pmod{13}$, откуда $\Delta_{13} \equiv 1 \pmod{13}$. Обращение в точный вычет локальной связности $13$ по макро-модулю $V=157$ ($2291554752 \equiv 13 \pmod{157}$) доказывает инвариант следа.
\end{proof}

\begin{law}[О мультипликаторах Холла и гомотопическом кручении]
Для ЦРМ $(157, 13, 1)$ любой простой делитель $q$ порядка $n_1 = k - \lambda = 12$ ($q \in \{2, 3\}$) является мультипликатором Холла. Минимальный делитель $q=3$ строго индуцирует свободную размерность коядра гомотопических препятствий $\Coker([\Phi_{10}] \ominus [\Phi_6]) \equiv 3 \pmod{13}$.
\end{law}

\subsection*{3. Спектральная теория ориентированных NoC}
Для ориентированного регулярного графа $\mathcal{G}$, индуцированного ЦРМ $(157, 13, 1)$ без учета тривиальной петли ($k_{\text{eff}} = 12$), спектральный радиус равен $\lambda_0 = 12.0$. Второе по модулю максимальное собственное значение составляет $\mu_{\text{max}} = 6.4270$. 

Все нетривиальные собственные значения лежат в комплексных границах жесткости ЦРМ. Алгебраическая связность направленной NoC-магистрали составляет:
\begin{equation}
\lambda_1 = k_{\text{eff}} - \mu_{\text{max}} = 12 - 6.4270 = 5.5730
\end{equation}
В соответствии с неравенством Чигера для ориентированных сетей, нижний порог пропускной способности сечений при деградации каналов удовлетворяет условию:
\begin{equation}
h(\mathcal{G}) \ge \frac{\lambda_1}{2} = \frac{5.5730}{2} = 2.7865
\end{equation}
Высокое значение $h(\mathcal{G})$ гарантирует сублогарифмическое перемешивание трафика и мгновенную перестройку путей обхода «на лету» без блокировок.

\subsection*{4. Научно-критический анализ}
\begin{criticism}[Адельное представление над составными базами]
Использование систем счисления с базами $B \in \{6, 12, 60\}$ в качестве адельных пространств является локальным изоморфизмом колец вычетов и схем над составными модулями $\mathbb{Z}/60\mathbb{Z}$, требующим осторожности при обобщении на полные группы аделей.
\end{criticism}

\begin{criticism}[Гомотопическая природа коядра]
Отождествление $\Coker = 3$ с кручением в группах Селмера выступает сильной геометрической моделью, требующей для полной строгости явного развертывания комплекса этальных когомологий, где разность операторов выполняет роль дифференциала.
\end{criticism}

\section*{Заключение}
Разработанная модель осуществляет переход от прикладной топологии NoC к основаниям математики. Аппарат проконечных пополнений и мотивов Гротендика доказывает, что аппаратная отказоустойчивость сетей на кристалле эквивалентна топологической жесткости ассоциированных арифметических схем.

\end{document}
